\chapter{Infraestrutura de Chave pública}
Um certificado digital possui o mesmo nível de confiança que a carteira de identidade. Sua principal função é garantir a autenticidade por meio da identificação digital de uma pessoa física ou jurídica. Duas entidades são necessárias para a validação de um certificado digital, a \textbf{autoridade de registro} e a \textbf{autoridade certificadora}. A primeira verifica a identidade do usuário que solicita o certificado digital, a segunda emite, armazena e assina os certificados digitais, funcionando como uma espécie de cartório digital. As autoridades certificadoras são organizadas hierarquicamente, com uma CA Raiz no topo delegando confiança para as CAs intermediárias.

Para se criar um certificado digital, o usuário gera localmente um par de chaves assimétricas, em seguida, gera uma \textbf{solicitação de assinatura de certificado (CSR)}, inserindo os dados de identidade e a chave pública. Esse pacote é resumido em um hash e assinado com a chave privada. Esse resumo criptográfico assinado é enviado para a CA que valida a assinatura, confirma a identidade por meio da RA e, por fim, assina o certificado final usando a chave privada da própria CA.

%\section{Algoritmos de Geração de Chaves}
%\subsection{Criptografia de Curvas Elípticas (ECC)}
%\subsection{Rivest-Shamir-Adleman (RSA)}
%\subsection{Diffie-Hellman}
\section{PKI baseada em Merkle Tree}
Diferentemente do RSA e ECC, os algoritmos de geração de chaves de certificados de merkle, por serem baseados em hash, possuem características pós-quântica. A eficiência matemática reside na forma como a validação ocorre. Se um usuário recebe um certificado e precisa provar que a folha específica desse certificado realmente faz parte da árvore original, ele não precisa baixar a árvore inteira. A AC envia junto do certificado o caminho de autenticação e os hashes dos nós irmãos desse caminho. Complexidade logarítmica ($O\log_2{N}$), com N sendo o número de folhas.

Podem ser aplicadas em esquemas stateful (LMS e XMSS) ou stateless (SLH-DSA e SPHINCS+). Os primeiros utilizam a árvore para agrupar um número finito de chaves de assinatura, cada nó é utilizado apenas uma vez, assim o sistema precisa armazenar quais nós já foram utilizados. Os segundos evitam esse problema gerando uma estrutura de árvores tão grande que a probabilidade de colisão seja quase nula.

\subsection{XMSS}

Possui base WOTS+, utiliza um contador de 0 até n folhas para evitar colisões. Antes de gerar o Hash, o algoritmo aplica máscaras de bit XOR sobre os nós filhos. $Hash((FilhoEsquerdo XOR Máscara1) || (FilhoDireito XOR Máscara2))$, assim, a função hash não precisa ser resistente a colisões, apenas resistente a segunda pre-imagem.

O pacote da assinatura XMSS possui três elementos:
\begin{enumerate}
	\item índice
	\item assinatura WOTS+ do documento
	\item Caminho de autenticação
\end{enumerate}

